\PassOptionsToPackage{dvipsnames,svgnames,table}{xcolor}
\documentclass[11pt]{article}
\usepackage[a4paper,margin=22mm]{geometry}
\usepackage[no-math]{fontspec}
\setmainfont{DejaVu Serif}
\usepackage{amsmath,amssymb,amsthm,mathtools,graphicx,xcolor,hyperref,xurl,xspace,ifthen,enumerate}
\usepackage[shortlabels]{enumitem}
\setlength{\emergencystretch}{4em}
\AtBeginDocument{\special{pdf:mapline pzdr Dingbats <uzdr.pfb}}
\SetMathAlphabet{\mathbf}{normal}{OT1}{cmr}{bx}{n}
\SetMathAlphabet{\mathbf}{bold}{OT1}{cmr}{bx}{n}
\newtheorem{theo}{Theorem}[section]
\newtheorem{lemm}[theo]{Lemma}
\newtheorem{thm}[theo]{Theorem}
\newtheorem{proposition}[theo]{Proposition}
\newtheorem{corollary}[theo]{Corollary}
\newtheorem{definition}[theo]{Definition}
\newtheorem{theorem}[theo]{Theorem}
\newtheorem*{remas*}{Remarks}
\newtheorem{remas}[theo]{Remarks}
\newtheorem*{exam*}{Example}
\newtheorem{ques}[theo]{Question}
\providecommand{\og}{«\nobreakspace}
\providecommand{\fg}{\nobreakspace»}
\newtheorem{lem}[theo]{Lemma}
\newtheorem{lemma}[theo]{Lemma}
\newtheorem{prop}[theo]{Proposition}
\newtheorem{coro}[theo]{Corollary}
\newtheorem{corol}[theo]{Corollary}
\newtheorem*{conj*}{Conjecture}
\newtheorem{conj}[theo]{Conjecture}
\newtheorem{conjecture}[theo]{Conjecture}
\newtheorem{Proposition}[theo]{Proposition}
\newtheorem{theoreme}[theo]{Theorem}
\theoremstyle{definition}
\newtheorem{defi}[theo]{Definition}
\newtheorem{exam}[theo]{Example}
\newtheorem{exem}[theo]{Example}
\newtheorem{remark}[theo]{Remark}
\newtheorem{example}[theo]{Example}
\newtheorem{exams}[theo]{Examples}
\newtheorem{rema}[theo]{Remark}
\newtheorem{nota}[theo]{Notation}
\newtheorem{notation}[theo]{Notation}
\newtheorem*{notas}{Notations}
\newtheorem{result}[theo]{Result}
\newtheorem{question}[theo]{Question}
\newtheorem{prob}[theo]{Problem}
\newtheorem*{rema*}{Remark}
\newtheorem*{defi*}{Definition}
\newtheorem*{theo*}{Theorem}
\newtheorem*{lemm*}{Lemma}
\newtheorem*{prop*}{Proposition}
\newcommand{\enoncename}{}
\newtheorem{corpusenonce}[theo]{\enoncename}
\newenvironment{enonce}[2][]{\renewcommand{\enoncename}{#2}\begin{corpusenonce}}{\end{corpusenonce}}
\newenvironment{enonce*}[2][]{\par\medskip\noindent\textbf{#2.}\itshape}{\par\medskip}
\newenvironment{proofc}[1][\proofname]{\begin{proof}[#1]}{\end{proof}}
\newcommand{\Mc}{\mathcal{M}}
\newcommand{\pointir}{.\ ---\ }
\newcommand{\equalenv}[2]{\expandafter\let\csname#1\expandafter\endcsname\csname#2\endcsname\expandafter\let\csname end#1\expandafter\endcsname\csname end#2\endcsname}
\providecommand{\diagbox}[3][]{\begin{tikzpicture}[x=1em,y=1em,baseline=(current bounding box.center)]\useasboundingbox(0,0) rectangle(3,1.4);\draw(0,1.4)--(3,0);\node at(.5,.3){#2};\node at(2.5,1.1){#3};\end{tikzpicture}}
\providecommand{\eg}{e.g.\ }
\providecommand{\ie}{i.e.\ }
\providecommand{\cf}{cf.\ }
\providecommand{\longthanks}[1]{\par\smallskip\noindent #1\par\smallskip}
\providecommand{\qbinom}[2]{\genfrac{[}{]}{0pt}{}{#1}{#2}}

\providecommand{\mathof}[1]{\mathrm{#1}}

\numberwithin{equation}{section}


\def\endappendix{}


\usepackage{latexsym,amsmath,amssymb,verbatim,tikz,tikz-cd}
\usepackage{mathtools}
\usepackage{graphicx,xcolor}

\usepackage[english]{babel}
\usepackage{amssymb}
\usepackage{amsmath}
\usepackage{amsthm}
\usepackage{graphicx}
\usepackage{csquotes}
\usepackage{bm}
\usepackage{enumitem}
\usepackage{subcaption}
\usepackage{caption}
\renewcommand{\thesubfigure}{\Alph{subfigure}}
\usepackage{dirtytalk}
\providecommand{\nopunct}{\spacefactor1007\relax}
\newtheorem{problem}[theo]{Problem}

\newcommand{\sym}{\mathrm{Sym}}
\newcommand{\qsym}{\mathrm{QSym}}
\newcommand{\Des}{\operatorname{Des}}
\newcommand{\Asc}{\operatorname{Asc}}
\newcommand{\Syt}{\operatorname{SYT}}
\newcommand{\Short}{\operatorname{Short}}
\newcommand{\Loop}{\operatorname{Loop}}
\newcommand{\Stable}{\operatorname{Stable}}
\newcommand{\len}{\operatorname{len}}
\newcommand{\height}{\operatorname{height}}
\newcommand{\core}{\operatorname{core}}
\newcommand{\res}{\operatorname{res}}
\newcommand{\row}{\operatorname{row}}
\newcommand{\insrt}{\operatorname{insert}}
\newcommand{\UP}{\operatorname{UP}}
\newcommand{\DOWN}{\operatorname{DOWN}}


\newcommand{\ZZ}{\mathbb{Z}}
\newcommand{\NN}{\mathbb{N}}
\newcommand{\RR}{\mathbb{R}}
\newcommand{\QQ}{\mathbb{Q}}
\newcommand{\PP}{\mathbb{P}}

\newcommand{\F}{\mathcal{F}}
\newcommand{\Q}{\mathcal{Q}}
\newcommand{\A}{\mathcal{A}}
\newcommand{\M}{\mathcal{M}}
\newcommand{\I}{\mathcal{I}}
\newcommand{\PM}{\mathcal{PM}}
\newcommand{\Pcal}{\mathcal{P}}


\newcommand{\odds}{\mathrm{odds}}

\makeatletter
\newcommand{\proofstep}[1]{
  \par
  \addvspace{\medskipamount}
  \textbf{#1\@addpunct{.}}\enspace\ignorespaces
}

\newcommand\restr[2]{{
  \left.\kern-\nulldelimiterspace 
  #1 
  \vphantom{\big|} 
  \right|_{#2} 
  }}
\makeatother












\title{The descent set of the superstandard tableau determines the dominance order of the shape}
\author{Avichai Marmor}
\date{doi:10.5802/alco.351; CC BY 4.0}
\begin{document}
\maketitle
\noindent\textit{Editorial note.} The counted statement is the author's labelled lemma `lem:superstandard characterize coefficients` (source lines 570-577); the proof environment begins after two blank lines, so the pairing is by adjacency.\par\medskip
\noindent\textit{Editorial note.} The retained context is the author's own lead-in sentence introducing the statement.\par\medskip
\noindent\textit{Editorial note.} The retained preamble is the editorially prepared standalone preamble retained for this source, which already carries the author's theorem-environment family and macro set; packages needed only by other parts of the article are left out because no retained line uses them.\par\medskip
\noindent\textit{Editorial note.} No retained line refers to a label, citation or figure outside the retained span, so the package declares no external reference.\par\medskip
\noindent\textit{Editorial scope.} One original result of the pinned publisher article, counted once: the characterisation of the superstandard tableau by its descent set (author label lem:superstandard characterize coefficients).\par\medskip
\noindent\textit{Allowed standard prerequisites (editorial).}\begin{itemize}
\item Superstandard tableau T\_lambda. T\_lambda is the standard Young tableau of shape lambda filled row by row with 1,2,...,N.. Both assertions are about T\_lambda and its descent set.
\item Descent set and conjugate partition. Des(T) is the set of positions i with i+1 in a lower row than i, and lambda' denotes the conjugate partition.. The statement identifies Des(T\_lambda) with [N-1] minus the partial sums of lambda', and the proof argues by columns of the conjugate shape.
\end{itemize}\medskip
The power of these notions may be reflected by the following statement:

\begin{lemma}
\label{lem:superstandard characterize coefficients}
    Let $\lambda, \mu \vdash N$ be two partitions. Then:
    \begin{enumerate}
        \item $\Des(T_\lambda) = [N-1] \setminus \{\lambda_{1}', \lambda_{1}'+\lambda_{2}', \dots\}$.
        \item If $\Des(T) = [N-1] \setminus \{\lambda_{1}', \lambda_{1}'+\lambda_{2}', \dots\}$ for some $T \in \Syt(\mu)$, then $\mu' \ge \lambda'$. Furthermore, if $\mu = \lambda$ then $T = T_\lambda$.
    \end{enumerate}
\end{lemma}

\begin{proof}
    The first assertion is obvious, so let us focus on the second assertion.
    
    Let $T \in \Syt(\mu)$ be a tableau, and assume that $\Des(T) = \Des(T_\lambda)$ and $T \ne T_\lambda$. We aim to show that $\mu' > \lambda'$. Let us denote the first column that differs between $T$ and $T_\lambda$ by $i$. The $i$-th column of $T_\lambda$ contains the entries $s+1,\dots,s+\lambda_{i}'$, where $s = \lambda_1' + \dots + \lambda_{i-1}'$. To prove $\mu' > \lambda'$, it suffices to show that all these entries also appear in the $i$-th column of $T$.
    
    Assume, by contradiction, that some of these entries appear in other columns of $T$. Let $x$ be the minimal such entry, and let $j \ne i$ be the column of $T$ that contains $x$. Since $s+1$ appears in the $i$-th column of $T$, we may assume that $x > s+1$. If $j < i$, this contradicts the assumption that $T$ and $T_\lambda$ agree in the first $i-1$ columns. On the other hand, if $j > i$, then $x$ appears in the first row of $T$, since it is the minimal entry of the $j$-th column of $T$. Consequently, we obtain that $x-1 \notin \Des(T)$, while $x-1 \in \Des(T_\lambda)$, which contradicts the assumption that $\Des(T) = \Des(T_\lambda)$.
\end{proof}

\end{document}
